
\documentclass{article}
\usepackage{colortbl}
\usepackage{makecell}
\usepackage{multirow}
\usepackage{supertabular}

\begin{document}

\newcounter{utterance}

\centering \large Interaction Transcript for game `cladder', experiment `full\_v1.5\_default', episode 9 with Qwen3.5{-}4B.
\vspace{24pt}

{ \footnotesize  \setcounter{utterance}{1}
\setlength{\tabcolsep}{0pt}
\begin{supertabular}{c@{$\;$}|p{.15\linewidth}@{}p{.15\linewidth}p{.15\linewidth}p{.15\linewidth}p{.15\linewidth}p{.15\linewidth}}
    \# & \multicolumn{2}{c}{Player} && \multicolumn{2}{c}{Game Master} \\
    \hline

    \theutterance \stepcounter{utterance}  
    & & & \multicolumn{4}{p{0.6\linewidth}}{
        \cellcolor[rgb]{0.9,0.9,0.9}{
            \makecell[{{p{\linewidth}}}]{
                \texttt{\tiny{[P1$\langle$GM]}}
                \texttt{You are an expert in causal inference. The following question is not a typical commonsense query, but rather a meticulously designed question created by a professor specializing in causal inference, intended to assess the students' mastery of the course content.} \\
\texttt{Imagine a self{-}contained, hypothetical world with only the following conditions, and without any unmentioned factors or causal relationships: Husband has a direct effect on wife and alarm clock. Wife has a direct effect on alarm clock. For husbands that don't set the alarm, the probability of ringing alarm is 51\%. For husbands that set the alarm, the probability of ringing alarm is 49\%. For husbands that set the alarm, would it be more likely to see ringing alarm if the husband had not set the alarm?} \\
\texttt{Start your answer with 'Yes' or 'No', followed by additional reasoning or evidence to support your explanation.} \\
            }
        }
    }
     \\ \\

    \theutterance \stepcounter{utterance}  
    & \multicolumn{4}{p{0.6\linewidth}}{
        \cellcolor[rgb]{0.561,0.737,0.561}{
            \makecell[{{p{\linewidth}}}]{
                \texttt{\tiny{[P1$\rangle$GM]}}
                \texttt{No} \\
\\ 
\texttt{**Reasoning:**} \\
\\ 
\texttt{This question is a classic counterfactual query designed to test the understanding of **selection bias** (specifically, Berkson's paradox or collider stratification) and the distinction between **causal effects** and **correlations**.} \\
\\ 
\texttt{1.  **Analyze the Causal Structure**:} \\
\texttt{    *   Let \$H\$ be the action of the husband (Set Alarm = 1, Don't Set = 0).} \\
\texttt{    *   Let \$A\$ be the state of the alarm (Rings = 1, Doesn't Ring = 0).} \\
\texttt{    *   The problem states:} \\
\texttt{        *   \$P(A=1 | H=0) = 0.51\$} \\
\texttt{        *   \$P(A=1 | H=1) = 0.49\$} \\
\texttt{    *   Crucially, the prompt establishes that the husband's action (\$H\$) is the *cause* of the alarm ringing (\$A\$). There is no third variable causing both \$H\$ and \$A\$. The relationship is purely causal: the decision to set the alarm influences the probability of it ringing.} \\
\\ 
\texttt{2.  **Evaluate the Counterfactual Question**:} \\
\texttt{    *   The question asks: "For husbands that set the alarm (\$H=1\$), would it be more likely to see ringing alarm if the husband had not set the alarm (\$H=0\$)?"} \\
\texttt{    *   Mathematically, this asks if \$P(A=1 | H=1) < P(A=1 | H=0)\$.} \\
\texttt{    *   From the given data:} \\
\texttt{        *   Probability of ringing given the alarm is set: **49\%**} \\
\texttt{        *   Probability of ringing given the alarm is not set: **51\%**} \\
\\ 
\texttt{3.  **Logical Deduction**:} \\
\texttt{    *   The question posits a hypothetical intervention on the *counterfactual* world where the husband *did not* set the alarm, while keeping the husband's identity/characteristics constant.} \\
\texttt{    *   In this counterfactual scenario (\$H=0\$), the probability of the alarm ringing is explicitly given as **51\%**.} \\
\texttt{    *   In the actual observed scenario (\$H=1\$), the probability of the alarm ringing is **49\%**.} \\
\texttt{    *   Comparing the two probabilities: \$51\textbackslash{}\% > 49\textbackslash{}\%\$.} \\
\texttt{    *   Therefore, in the counterfactual world where the husband did not set the alarm, the alarm is **more likely** to ring than in the actual world where he did set it.} \\
\\ 
\texttt{4.  **Addressing Potential Confusion (The "Yes" Trap)**:} \\
\texttt{    *   A student might intuitively think, "If he didn't set the alarm, it wouldn't ring," leading them to answer "No." This ignores the specific conditional probabilities provided in the prompt} \\
            }
        }
    }
    & & \\ \\

    \theutterance \stepcounter{utterance}  
    & & & \multicolumn{2}{p{0.3\linewidth}}{
        \cellcolor[rgb]{0.9,0.9,0.9}{
            \makecell[{{p{\linewidth}}}]{
                \texttt{\tiny{[GM$|$GM]}}
                \texttt{No} \\
            }
        }
    }
    & & \\ \\

    \theutterance \stepcounter{utterance}  
    & & & \multicolumn{2}{p{0.3\linewidth}}{
        \cellcolor[rgb]{0.9,0.9,0.9}{
            \makecell[{{p{\linewidth}}}]{
                \texttt{\tiny{[GM$|$GM]}}
                \texttt{yes} \\
            }
        }
    }
    & & \\ \\

    \theutterance \stepcounter{utterance}  
    & & & \multicolumn{2}{p{0.3\linewidth}}{
        \cellcolor[rgb]{0.9,0.9,0.9}{
            \makecell[{{p{\linewidth}}}]{
                \texttt{\tiny{[GM$|$GM]}}
                \texttt{game\_result = LOSE} \\
            }
        }
    }
    & & \\ \\

\end{supertabular}
}

\end{document}
